\begin{table}[htbp]
\centering
\caption{Replacement Quality by Relative Draftee Wage}
\label{tab:replacement-quality-relative-draftee-wage}
\scriptsize
\begin{threeparttable}
\resizebox{\textwidth}{!}{%
\begin{tabular}{lcccc}
\toprule
 & Baseline wage percentile & Baseline salary & Baseline position rank & Baseline court rank \\
 & (1) & (2) & (3) & (4) \\
\midrule
Low relative-wage draftee destination & -0.1423$$ & -7.4901$$ & 0.0449$$ & 0.0870$$ \\
                                      & (0.0909) & (12.0346) & (0.0522) & (0.1546) \\[3pt]
Female &  &  & -0.4266$^{***}$ & -0.1535$$ \\
       &  &  & (0.1542) & (0.1019) \\[3pt]
Tenure (yrs) & -0.0031$$ & -0.1247$$ & 0.0263$^{***}$ & -0.0015$$ \\
             & (0.0162) & (2.1198) & (0.0065) & (0.0153) \\[3pt]
Different kyoku transfer & -0.1435$$ & -7.5658$$ & 0.0209$$ & 0.1493$^{**}$ \\
                         & (0.1248) & (7.8690) & (0.0786) & (0.0753) \\[3pt]
\midrule
Observations & 29 & 29 & 200 & 200 \\
$R^2$ & 0.146 & 0.036 & 0.168 & 0.006 \\
\bottomrule
\end{tabular}
}%
\begin{tablenotes}[flushleft]\footnotesize
\item \textit{Notes:} Sample consists of non-drafted workers who transferred to a draft-vacancy destination office-year. Draftee wage is measured as the drafted worker's wage percentile within occupation-by-year cells, then averaged across drafted workers in the destination office-year. The omitted category is a high relative-wage draftee destination, defined as at-or-above the median destination office-year draftee wage percentile. Heteroskedasticity-robust standard errors are in parentheses. $^{***}p<0.01$, $^{**}p<0.05$, $^{*}p<0.1$. Outcomes measure the transferee's pre-transfer quality using lagged/baseline characteristics. Baseline wage percentile is measured within the transferee's own baseline occupation-year. Baseline salary columns use only workers with observed lagged salary.
\end{tablenotes}\end{threeparttable}\end{table}
